% Separate unassembled candidate for the multicomponent count, not a completed count.
\subsection{Brauer classes and residual blocks on an exceptional tree}

\begin{lemma}[Proved: tree residues and the actual coefficient gerbe]
\label{D:SNC-tree-Brauer-geometry}
Let $R$ be a complete excellent normal local surface over an
algebraically closed coefficient and residue field $k$ of
characteristic $p>0$. Let $Y\to\operatorname{Spec}R$ be a
projective regular resolution, an isomorphism over the puncture
$U$, and let $E=\bigcup_{i=1}^v E_i$ be its reduced exceptional
divisor. Suppose $E$ is strict normal crossings and its dual
graph is a tree. Write $g_i=g(E_i)$ and
$N=(-(E_i\cdot E_j))_{ij}$.

For every $n$ prime to $p$, residue gives an isomorphism
\[
 \operatorname{Br}'(U)[n]\simeq
       \bigoplus_i H^1(E_i,\mathbb Z/n),\qquad
 |H^2(U,\mu_n)|=|\operatorname{coker}(N\bmod n)|
                             n^{2\sum_i g_i}.           \tag{D.TR.1}
\]
In particular the latter order is at most
$|\det N|n^{2\sum_i g_i}$. All secondary residues at the
nodes vanish.

A class $\beta$ of prime-to-$p$ order $m$ is killed by a
connected cyclic degree-$m$ cover $Y'\to Y$ which is etale
on all of $Y$. Its restriction to $U$ is the puncture of a
finite complete local normal extension of $R$. The exceptional
divisor upstairs need not have a tree as its dual graph.

Let $\alpha\in H^2(U,\mu_n)$ have Brauer image $\beta$,
and let $m_i$ be the order of the $i$th residue, with
$m_i=1$ for zero residue. Then $m_i\mid m\mid n$, and
the actual class $\alpha$ extends to
\[
       \mathcal X=\sqrt[m_1]{(Y,E_1)}\times_Y\cdots
                         \times_Y\sqrt[m_v]{(Y,E_v)}.  \tag{D.TR.2}
\]
At a node $E_i\cap E_j$, the restriction of its gerbe to
$B(\mu_{m_i}\times\mu_{m_j})$ has zero geometric Brauer
class. Its inertia is therefore a commutative diagonalizable
extension of $\mu_{m_i}\times\mu_{m_j}$ by $\mu_n$.
No nontrivial node commutator is discarded in this statement:
the tree hypothesis forces that commutator to vanish.
\end{lemma}

\begin{proof}
Put $\Lambda=\mathbb Z/n$. Normalization of the proper
nodal curve $E$ gives
\[
 0\longrightarrow\Lambda_E\longrightarrow
   \bigoplus_i\Lambda_{E_i}\longrightarrow
   \bigoplus_{x\in\operatorname{Nodes}(E)}\Lambda_x
   \longrightarrow0.                                  \tag{D.TR.3}
\]
The last map is the difference of the two branches, for
any fixed orientations. Proper strict-local base change,
nilpotent invariance, and the Picard lifting argument of
Lemma~\ref{D:local-Picard-prime-to-p-generators} show
\[
 H^2(Y,\mu_n)\simeq H^2(E,\mu_n)\simeq\Lambda^v,
 \qquad \operatorname{Pic}(Y)/n\simeq\Lambda^v,
 \qquad H^3(Y,\mu_n)=0.                               \tag{D.TR.4}
\]
Here multidegree is surjective on $\operatorname{Pic}(E)$:
choose a smooth point on each component. Its degree-zero
kernel is an extension of the product of the Jacobians by
the graph torus, hence is $n$-divisible. The formal Picard
kernel is uniquely $n$-divisible. In particular
$\operatorname{Br}'(Y)[n]=0$. These arguments impose no
finite-type-over-$k$ hypothesis on the complete local ring.

Apply derived Hom into $\mu_n$ on $Y$ to (D.TR.3),
viewing its closed-stratum sheaves as direct images to $Y$.
Closed-immersion adjunction identifies the resulting
triangle with the Mayer--Vietoris triangle for supports.
Absolute purity gives shifts $[-2]$ for a component and
$[-4]$ for a node. Thus its degree-three exact sequence is
\[
 0\longrightarrow\bigoplus_i H^1(E_i,\Lambda)
 \longrightarrow H^3_E(Y,\mu_n)
 \longrightarrow\Lambda(-1)^{e}
       \xrightarrow{\partial}\Lambda(-1)^v,            \tag{D.TR.5}
\]
where $e$ is the number of nodes. The map $\partial$ is
the signed incidence map: purity is compatible with
composition of the point-to-curve and curve-to-surface
Gysin maps, and the point class has degree one on each
proper component. This uses the exact regular-scheme
scope and compatibility of \citet[Expose XVI, Sections2--3,
Theorem3.1.1 and the proof of Corollary3.1.3]{A_ILO2014}.
The closed-stratum exact sequence used in that proof is
also (D.TR.3). All Tate twists are constant over $k$.

For a tree, $\partial$ is injective over every $\Lambda$.
Indeed a leaf row has just one nonzero entry, equal to
$1$ or $-1$; remove that edge and repeat. Hence
$H^3_E(Y,\mu_n)=\bigoplus_i H^1(E_i,\Lambda)$.
This calculation also explains the secondary-residue
assertion: the node values in (D.TR.5) form a balanced
flow, and a tree has no nonzero such flow.

Every line on the regular puncture extends to $Y$ by
closing a Cartier divisor. Consequently localization
and (D.TR.4) identify
\[
 \operatorname{Br}'(U)[n]
 =H^2(U,\mu_n)/(\operatorname{Pic}(U)/n)
                         =H^3_E(Y,\mu_n).              \tag{D.TR.6}
\]
The kernel of $\operatorname{Pic}(Y)\to\operatorname{Pic}(U)$
is generated by the exceptional divisor lines. Under the
multidegree identification in (D.TR.4), their map is the
intersection matrix $-N$. Therefore
$\operatorname{Pic}(U)/n=\operatorname{coker}(N\bmod n)$.
Smith normal form bounds its order by $|\det N|$;
negative definiteness makes this determinant nonzero.
This proves (D.TR.1). In addition, one fixed finite list
of line bundles on $Y$ represents every
$\operatorname{Pic}(U)/n$: choose representatives of
the finite group $\operatorname{coker}N$ once and use
right exactness of tensoring with $\mathbb Z/n$.

The same normalization sequence gives
$H^1(E,\mathbb Z/m)\simeq\bigoplus_iH^1(E_i,\mathbb Z/m)$
when the graph is a tree. Thus the residue tuple of a
period-$m$ class is a character on $E$ of exact order
$m$. Its torsor is connected and cyclic of degree $m$.
The proper henselian equivalence for finite etale covers
lifts it, including its action, to $Y'\to Y$
\citep[Lemma58.9.1, Tag0A48]{StacksSingleCurveBrauerD}.
Properness and connectedness of the closed fibre show
$Y'$ is connected; regularity then shows it is integral.
Its Stein ring is finite, complete, local and normal,
as in Proposition~\ref{D:single-curve-Brauer-residue-cover}.

The pulled-back class has zero generic residue on every
exceptional component upstairs. It is unramified elsewhere
because it was defined on the puncture. Purity on the
regular surface therefore extends it to $Y'$, whose
prime-to-$p$ Brauer group is zero by the same Picard and
proper-base-change argument as (D.TR.4). This proves
splitting without asserting that the new dual graph is
a tree. Indeed a connected etale cover of a nodal tree
of positive-genus curves can acquire graph cycles.

At the generic point of $E_i$, taking the $m_i$th root
of its parameter kills exactly its cyclic residue.
The proof of \citet[Proposition3.2.1, complete pp10--11]
{Lieblich2011ResidualD} applies with these componentwise
orders: the DVR construction is local on each component,
and regular-stack purity fills the remaining codimension-two
set. It gives a unique Brauer extension on (D.TR.2).
Lift it by Kummer to $H^2(\mathcal X,\mu_n)$ and correct
its restriction by a line from $\operatorname{Pic}(Y)$.
This is possible by the just-proved surjection onto
$\operatorname{Pic}(U)$, and extends the actual $\alpha$.

At a node, restrict to the strictly henselian regular
local surface. Each component residue is etale at the
node and becomes zero on that strictly local branch.
All divisor residues therefore vanish. Purity extends
the local Brauer class over the regular local surface,
where it is zero: an Azumaya algebra splits over a
strictly henselian local ring. Injectivity for the
regular root stack shows its node restriction also has
zero Brauer image. Kummer then says its class on
$B(\mu_{m_i}\times\mu_{m_j})$ is the root class of a
character line. The inertia of that root gerbe is
\[
 \{(h,z): h\in\mu_{m_i}\times\mu_{m_j},\quad
                       z^n=\lambda(h)\},              \tag{D.TR.7}
\]
for that character $\lambda$. This is commutative and
diagonalizable; all its orders are prime to $p$.
It proves the final assertion including the possibly
nonzero scalar Kummer part at the node.
\end{proof}

\begin{lemma}[Proved: the residual tree category and its node gluing]
\label{D:SNC-tree-residual-category}
In the preceding setting let $\mathcal G\to\mathcal X$
be the actual coefficient gerbe. On the normalization
of its reduced exceptional curve, the scalar-weight-one
category on component $i$ is equivalent to vector bundles
on
\[
 \mathcal R_i=\sqrt[m_j]{(R_i,P_{ij})}_{j\sim i},      \tag{D.TR.8}
\]
after choice of a neutralizing line. Here $R_i\to E_i$
is the connected etale cover of degree $m_i$ determined
by the residue, and $P_{ij}$ is the set of its $m_i$
points above the node with $E_j$. For multiple edges the
points are counted separately; the tree assumption has
already excluded parallel edges. The ordinary spectral
rank is $s_i=r/m_i$ for a rank-$r$ scalar-weight-one
bundle, so $m_i\mid r$.

Bundles on the reduced exceptional gerbe are precisely
these component bundles together with equivariant
isomorphisms of their fibres at every node. For fixed
joint character multiplicities $r_{ab}$ at a node,
after choosing homogeneous frames for the character
fibres, the isomorphisms and their inverses have polynomial
witnesses with at most $2\sum_{a,b}r_{ab}^2\leq2r^2$
coordinates and quadratic equations. This frame-atlas
description is valid over arbitrary coefficient algebras;
without frames, the character fibres may be projective
modules. It is an actual
reduced-curve equivalence; it asserts no realization
of normal deformations or final representation count.
\end{lemma}

\begin{proof}
The root-$m_i$ spectral construction on the generic
$E_i$ extends over its whole root-marked curve. Its
character torsor has no monodromy from an adjacent
root inertia: that monodromy is precisely the
commutator in (D.TR.7), which is zero. It is therefore
the pullback of $R_i\to E_i$. Applying the exact
character decomposition of
Proposition~\ref{D:residual-twisted-spectral-category}
on this base yields a Picard gerbe over (D.TR.8).
Its geometric Brauer group is zero. Indeed a root
of a smooth point on a smooth proper curve over $k$
has relative $R^1\pi_*\mathbf G_m=\mathbb Z/m_j$
supported at that point and $R^2\pi_*\mathbf G_m=0$;
the latter is the algebraically closed instance of
\citet[Lemmas4.1.1--4.1.3 and Proposition4.1.4]
{Lieblich2011ResidualD}. The base curve has zero
Brauer group and the point sheaves have zero first
cohomology. The Leray sequence therefore gives
$\operatorname{Br}'(\mathcal R_i)=0$ prime to $p$.
Neutralizing that Picard gerbe gives the claimed
category. Its inverse is pushforward of the universal
weight line tensored with the spectral bundle, exactly
as in the already proved spectral proposition.
The rank formula is $r=m_i s_i$.

At a node all the inertia is the fixed diagonalizable
group (D.TR.7). Scalar-weight-one modules split into
its characters over every $k$-algebra, by the exact
grading projections. Along either branch their
restrictions give the same joint-character list
exactly when the corresponding multiplicities agree.
After passing to the homogeneous fibre-frame atlas,
for a given joint character $(a,b)$ matching fibres
is a matrix of size $r_{ab}$ and an inverse matrix,
with the two product identities. Summing the blocks
gives the asserted coordinate and degree bounds.
The character labels include the fixed central
extension (D.TR.7); they are not replaced by an
unmarked pair of cyclic groups.

Finally the polynomial model of an ordinary node has
the fibre-product description
\[
       A[x,y]/(xy)=A[x]\times_A A[y]                 \tag{D.TR.9}
\]
for every coefficient algebra $A$, and the same identity
holds on compatible localizations. For branch modules
with a specified fibre identification, the gluing module
is the kernel of their difference map to the identified
node fibre. Locally on the coefficient and branch frame
atlases, change one basis by the prescribed invertible
matching matrix; this kernel is then a free module over
the fibre product. This is the elementary Milnor pasting
construction, with inverse given by restriction to the
two branches. The same proof is equivariant for the
diagonalizable finite node group, by choosing homogeneous
frames. These constructions commute with coefficient
changes. Applying them on polynomial node charts and
their etale localizations, followed by etale descent,
then gives vector bundles on the reduced nodal
exceptional gerbe, proving the asserted equivalence.
It does not replace the two-dimensional normal
neighbourhood by its reduced node.
\end{proof}

\begin{remark}[The remaining quantitative step]
On $\mathcal R_i$, normalized root-point weights encode
ordinary rank-$s_i$ bundles and flags at $O(m_i)$ marked
points. Their Hom sheaf is the kernel of the ordinary
Hom sheaf mapping to the forbidden-weight blocks at
those points. Thus its cohomology is represented by a
fixed-width cone of ordinary curve Hom and finite point
evaluations; arbitrary flag-dependent inverse matrices
are not automatically permitted in a quadratic bound.

To finish a tree count one must still construct actual
all-ring normal bundle families on a fixed ordinary
exceptional cycle and prove uniform slope windows and
Hom tails for this residual category. The mixed-root
numerical budget alone is adequate: if $m_i\mid m$,
weight the $i$th component's layer-$j$ variables by
$E(m/m_i)j$ and node matrices by $Em$. With common
equation degree $O(Em)$, each component has the same
factorial cancellation as before, and node variables
cost a constant per coordinate. This observation
does not prove the necessary weighted equations,
normal realization, or a tree representation count.
\end{remark}
